\section{Higher-Dimensional Torus Translations: Exact Closure Dimension Formula and Number-Field Constraints}

\subsection{Closure Subtori and Dimension Formula}

On $\TT^d=\RR^d/\ZZ^d$ consider the translation
\[
T_{\boldsymbol{\alpha}}(\mathbf{x})=\mathbf{x}+\boldsymbol{\alpha}\pmod 1,\qquad \boldsymbol{\alpha}=(\alpha_1,\dots,\alpha_d)\in\RR^d.
\]
Define the annihilator lattice
\[
L(\boldsymbol{\alpha})=\{\mathbf{m}\in\ZZ^d:\ \mathbf{m}\cdot\boldsymbol{\alpha}\in\ZZ\}.
\]
By character annihilation theory, the orbit closure is a subtorus
\[
H=\{\mathbf{x}\in\TT^d:\ \mathbf{m}\cdot \mathbf{x}=0\ \text{in}\ \TT\ \text{for all}\ \mathbf{m}\in L(\boldsymbol{\alpha})\},
\]
with
\[
\dim H=d-\mathrm{rank}\,L(\boldsymbol{\alpha}).
\]

To obtain an expression directly related to algebraic structure, introduce
\[
r(\boldsymbol{\alpha})=\dim_{\QQ}\mathrm{span}_{\QQ}\{1,\alpha_1,\dots,\alpha_d\}.
\]

\begin{theorem}[Closure Dimension and $\QQ$-Rank]\label{thm:closure-dim}
For any $\boldsymbol{\alpha}\in\RR^d$, the dimension of the orbit closure subtorus $H$ satisfies
\[
\dim H=r(\boldsymbol{\alpha})-1.
\]
\end{theorem}
See Appendix \ref{app:closure-dim-proof} for the proof.

\subsection{Corollaries from Number-Field Constraints}

Let $K$ be a number field of degree $k$. If $\alpha_1,\dots,\alpha_d\in K$, then the $\QQ$-linear dimension of $\{1,\alpha_1,\dots,\alpha_d\}\subset K$ is at most $k$, i.e., $r(\boldsymbol{\alpha})\le k$, hence:

\begin{corollary}[Number-Field Degree Upper Bound]
If $\alpha_i\in K$ and $[K:\QQ]=k$, then
\[
\dim H\le k-1.
\]
\end{corollary}

In particular, when $k=2$ (quadratic fields), the closure dimension is at most 1, meaning the orbit closure is at most a one-dimensional subtorus. The golden field $K=\QQ(\sqrt 5)$ falls into this case.

\subsection{Restrictiveness of Golden Power Coordinates}

If we take $\boldsymbol{\alpha}=(\alpha,\alpha^2,\dots,\alpha^d)$ with $\alpha=\varphi^{-1}$, since $\alpha$ satisfies the quadratic relation $\alpha^2+\alpha=1$, the $\QQ$-linear dimension of $\{1,\alpha,\alpha^2,\dots,\alpha^d\}$ is 2. By Theorem \ref{thm:closure-dim} we obtain
\[
\dim H=1.
\]
Thus this type of ``power coordinates within a single quadratic field'' cannot generate a full-rank dense orbit on $\TT^d$; its coverage dimension is strictly limited. This constraint provides a verifiable bound on ``the filling capacity of single-step scanning in multi-dimensional space''.

\subsection{Typical Growth Rates and Comparison with Algebraic Special Points}

For Lebesgue-almost-every $\theta$, the exponential growth rate of convergence denominators $q_n(\theta)$ exists and is a constant
\[
\lim_{n\to\infty}\frac{1}{n}\log q_n(\theta)=\frac{\pi^2}{12\log 2},
\]
hence
\[
\lim_{n\to\infty} q_n(\theta)^{1/n}=\exp\!\left(\frac{\pi^2}{12\log 2}\right).
\]
For quadratic irrationals, the $q_n$ growth is determined by the spectral radius of the period matrix. This comparison serves to distinguish ``statistical renormalization of typical parameters'' from ``strict self-similarity of algebraic periodic points''.
